\documentclass[12pt,a4paper]{exam}

% --------------------
% Packages
% --------------------
\usepackage[utf8]{inputenc}
\usepackage[T1]{fontenc}
\usepackage[french]{babel}
\usepackage{amsmath,amssymb}
\usepackage{mdframed}
\usepackage{geometry}
\usepackage{pgfplots}
\usepackage{tkz-tab}
\usepackage{array}
\pgfplotsset{compat=1.18}
\geometry{margin=2cm}

% Pour les coordonnées de vecteurs :
\newcommand*{\Coord}[2]{% 
  \ensuremath{
    \begin{pmatrix} 
      #1\\[5pt]
      #2 
    \end{pmatrix}}}

% -------------------
% En-tête / pied (exam)
% --------------------
\header{}{}{2\up{nde}10}
\footer{}{\thepage}{}
\setlength{\parindent}{0pt}

% --------------------
\begin{document}
\begin{mdframed}[
    linewidth=0.5pt,
    roundcorner=5pt,
    % linecolor=ForestGreen
]
\begin{center}
    \Large\textbf{Contrôle sur les vecteurs\\2nde générale}
\end{center}

\vspace{0.3cm}

\noindent
\textbf{Durée :} 1h \hfill
\textbf{Calculatrice :} \underline{NON} autorisée
\end{mdframed}

\vspace{20pt}

\textbf{\underline{Exercice 1} (9 points)}

\begin{center}
    \begin{tikzpicture}[scale=1]
    % Axes du repère
    \draw[->, thick] (-7,0) -- (7,0) node[right]{$x$};
    \draw[->, thick] (0,-3) -- (0,4) node[above]{$y$};

    % Origine
    \node at (0,0) [below left]{$O$};

    % Vecteurs du repère
    \draw[->, blue, very thick] (0,0) -- (1,0) node[below left]{$\overrightarrow{i}$};
    \draw[->, red, very thick] (0,0) -- (0,1) node[below left]{$\overrightarrow{j}$};

    % Vecteurs personnalisés
    \draw[->, black, very thick] (-3,3) -- (-5,1) node[pos=0.5, above left]{$\overrightarrow{u}$};
    \draw[->, black, very thick] (1,3) -- (3,3) node[pos=0.5, above]{$\overrightarrow{v}$};
    \draw[->, black, very thick] (1,-1) -- (4,0) node[pos=0.5, above left]{$\overrightarrow{w}$};

    % Grille (optionnelle)
    \draw[step=1cm, gray, very thin] (-7,-3) grid (7,4);
    \end{tikzpicture}
\end{center}

\vspace{5pt}
\begin{questions}

\question[0,5] Placer les points $A\Coord{2}{1}$ et $B\Coord{-1}{2}$ sur le graphique.
\begin{parts}
    \part[0,5] Caluler les coordonnées de $\overrightarrow{AB}$.
    \part[1] Calculer la distance $AB$.
\end{parts}

\vspace{10pt}

\question
\begin{parts}
    \part[2] Lire graphiquement les coordonnées des vecteurs $\overrightarrow{u}$, $\overrightarrow{v}$ et $\overrightarrow{w}$.
    \part[1] Calculer les coordonnées de $2\overrightarrow{u} - 3\overrightarrow{v}$.
\end{parts}

\question[1] Représenter \textbf{sur le graphique} le vecteur : $\overrightarrow{k}\Coord{1}{1}$

\question[2] Les vecteurs $\overrightarrow{u}$ et $\overrightarrow{k}$ sont-ils colinéaires ? 
\underline{Justifier.}

\question[1] Calculer la norme du vecteur $\overrightarrow{w}$.



\end{questions}

\vspace{1cm}

% ========================================================
%             Exercice 2
% ========================================================

\textbf{\underline{Exercice 2} (9 points)}

On considère les points $A\Coord{-2}{4}$, $B\Coord{0}{6}$ et $C\Coord{4}{6}$.

\begin{questions}
    \question[2] Déterminer les coordonnées du points $D$ tel que $ABCD$ est un parallélogramme.
    \question[2] On pose le point $E\Coord{6}{4}$. Montrer que le droites $(CE)$ et $(BD)$ sont parallèles.
    \question[2] Montrer que $CED$ est un triangle isocèle en $C$.
    \question[1] Calculer les coordonnées du point $I$ milieu de $[BE]$.
    \question[2] Les points $D$, $I$ et $C$ sont-ils alignés ? \underline{Justifier}.
\end{questions}

\vspace{1cm}

% ========================================================
%             Exercice 3
% ========================================================
\noindent
\textbf{\underline{Exercice 3} (2 points)}

\vspace{5pt}
Soient $\vec{m}$ et $\overrightarrow{n}$ deux vecteurs du plan.

\begin{questions}

\question[2] Les vecteurs $-12\overrightarrow{m}+4\overrightarrow{n}$ et $9\overrightarrow{m}-3\overrightarrow{n}$ sont-ils 
colinéaires ?

\end{questions}

\vspace{1cm}

% ========================================================
%             BONUS
% ========================================================
\textbf{\underline{Exercice BONUS} (2 points)}

$ABCD$ est un parallélogramme de centre $O$. $E$ est le point tel que $\overrightarrow{AE} = 3 
\overrightarrow{AB}$ et $F$ le point tel que $\overrightarrow{CF} = -2 \overrightarrow{AB} - 
\dfrac{1}{5}\overrightarrow{AD}$.

\begin{questions}
    \question[1] Démontrer, en utilisant la relation de Chasles, que $\overrightarrow{FE} = 
    4 \overrightarrow{AB} - \dfrac{4}{5}\overrightarrow{AD}$.

    \question[1] Démontrer de même que $\overrightarrow{FO} = \dfrac{3}{2}\overrightarrow{AB} - 
    \dfrac{3}{10}\overrightarrow{AD}$.

    \question[1] En déduire que les points $F$, $O$ et $E$ sont alignés.
\end{questions}




\vspace{0.5cm}

\begin{center}
    \textit{Fin de l'épreuve}
\end{center}

\end{document}
